filtered @@M0@@-algebras @@M1@@ by sending @@M2@@ to @@M3@@ and @@M4@@ to the @@M5@@-th coefficient of @@M6@@ for @@M7@@; using this homomorphism, we see that it suffices to verify that the morphism of functors under consideration vanishes on the element @@M8@@. We then use the endomorphism @@M9@@ of @@M10@@ defined by @@M11@@ for every @@M12@@, and the commutative diagram:
@@M13@@
now @@M14@@, where @@M15@@ is the augmentation homomorphism of @@M16@@; since the latter annihilates @@M17@@, the analogous diagram written for @@M18@@ shows that @@M19@@ annihilates the image of @@M20@@ in @@M21@@, giving the result.

Passing to the quotient, we therefore associate to every morphism of functors: @@M22@@ a morphism of functors: @@M23@@; by composing with @@M24@@, and using the isomorphism @@M25@@, we obtain a morphism: @@M26@@, which defines the required homomorphism @@M27@@. It is immediate to verify that @@M28@@ and @@M29@@ are inverses of one another.

\sgaStableTarget{sga6.e05.thm.1-13}\textbf{Theorem 1.13.} \uline{Let} @@M30@@ \uline{be a functor on} @@M31@@, \uline{with values in the category of separated complete filtered abelian groups}, \uline{and commuting with invariants under a finite group}. \uline{Then the map from} @@M32@@ \uline{to} @@M33@@ \uline{which}, \uline{to a morphism of functors}, \uline{assigns its value on the element}
